% ===== CHARACTER REGISTER =====
\newglossaryentry{Kuonrat von Steinare}{
  type=characters,
  name={Kuonr\^{a}t von Stein\ae{}re},
  description={Ritter aus dem Geschlecht der Ottensteiner, geboren um 1213 auf Burg Ottenstein im Waldviertel. Sohn Hadmars von Stein\ae{}re (~1180--~1238), eines jüngeren Sohnes des urkundlich belegten Burgherrn Hugo de Ottenstaine (1177/92); obwohl erstgeborener Sohn nicht Erbe der Burg. Selbsternannter Streiter der \textit{vr\^{\i}heit}. \textbf{Fiktionale Figur (historisches Umfeld zulässig)}}
}

\newglossaryentry{Hadmar von Steinare}{
  type=characters,
  name={Hadmar von Stein\ae{}re},
  description={Ritter, jüngerer Sohn des urkundlich belegten Burgherrn Hugo de Ottenstaine (1177/92); kein Erbanspruch auf Burg Ottenstein. Geboren um 1180, gestorben um 1238. Vater \gls{Kuonrat von Steinare}s. \textbf{Fiktionale Figur (historisches Umfeld zulässig)}}
}

\newglossaryentry{Veltloufer}{
  type=characters,
  name={Veltl\^{o}ufer},
  description={Das Streitross Kuonr\^{a}ts von Stein\ae{}re. Sein Name \textit{veltl\^{o}ufer} bedeutet der \enquote{Feldläufer}. \textbf{Fiktionale Figur}}
}

\newglossaryentry{Hugo von Ottenstaine}{
  type=characters,
  name={Hugo von Ottenstaine},
  description={Burgherr zu Ottenstein im Waldviertel, urkundlich belegt 1177 und 1192. Großvater \gls{Kuonrat von Steinare}s. \textbf{Historisch belegte Figur}}
}

\newglossaryentry{Wolfhart von Raabs}{
  type=characters,
  name={Wolfhart von Raabs},
  description={Niederadeliger Grundherr zu Raabs an der Thaya im Waldviertel. Vater \gls{Denise von Raabs}s. \textbf{Fiktionale Figur (historisches Umfeld zulässig)}}
}

\newglossaryentry{Denise von Raabs}{
  type=characters,
  name={Den\^{\i}se von Raabs},
  description={Niedere Adelige, Tochter \gls{Wolfhart von Raabs}s. Einstige Verlobte \gls{Kuonrat von Steinare}s. \textbf{Fiktionale Figur}}
}

\newglossaryentry{Konrad von Hardegg}{
  type=characters,
  name={Konrad von Hardegg},
  description={Graf von Plain und Hardegg, wohl um 1180 geboren und am 4. April 1250 verstorben. Jüngerer Sohn Heinrichs I. von Plain-Hardegg (Haus der Wilhelminer), urkundlich als Graf von Plain und Hardegg zwischen dem spaeten 12. und der Mitte des 13. Jahrhunderts belegt. Für die Romanfassung als streitsüchtige, dem Trunk verfallene und frauenverachtende Figur mit triebhaftem Verhalten zugespitzt. \textbf{Historische Figur}}
}

\newglossaryentry{Dirc Kruter}{
  type=characters,
  name={Dirc Kr\^{u}ter},
  description={Fahrender Quacksalber, der behauptet, Stein in Gold verwandeln zu können. \textbf{Fiktionale Figur}}
}

\newglossaryentry{Kuono}{
  type=characters,
  name={Kuono},
  description={Junger Weltgeistlicher, intelligent, von freiem Willen Priester geworden. Hadert mit der Keuschheit. Gefährte \gls{Kuonrat von Steinare}s auf dessen Reisen. \textbf{Fiktionale Figur}}
}

\newglossaryentry{Einsidel}{
  type=characters,
  name={Der Einsidel},
  description={Namenloser alter Büßer, der in einer Felsklause im Thayatal nahe Hardegg lebt. Ehemaliger Ritter aus der Steiermark; erschlug in blinder Eifersucht seinen eigenen Sohn, den er mit dem Liebhaber seiner Frau verwechselt hatte. Angelehnt an die lokale Einsiedlersage der Einsiedlerklause im Thayatal (realer Ort, urkundlich bewohnt). \textbf{Sagengestalt}}
}

\newglossaryentry{Nikolaus von Steinare}{
  type=characters,
  name={Nicl\^{a}s von Stein\ae{}re},
  description={Jüngerer Bruder \gls{Kuonrat von Steinare}s; Erbe der väterlichen Güter des Hauses von Stein\ae{}re. \textbf{Fiktionale Figur (historisches Umfeld zulässig)}}
}

\newglossaryentry{Gruober}{
  type=characters,
  name={Gruober},
  description={Fahrender Spielmann, der zwischen Böhmen und Österreich zieht. Entführt Frauen in Böhmen und verkauft sie als Prostituierte. \textbf{Fiktionale Figur}}
}

\newglossaryentry{Knecht}{
  type=characters,
  name={Knecht},
  description={Namenloser Diener im Gefolge \gls{Kuonrat von Steinare}, von diesem ausschließlich als \enquote{Knecht} bezeichnet. Unterwürfiger Charakter, der von seinem Herrn schlecht und herabsetzend behandelt wird. \textbf{Fiktionale Figur}}
}
